% !TEX root = ../main.tex
% Related Work
% Claims to support: RAG for classification, ML4SE issue triage, retrieval debiasing

\section{Related Work}
\label{sec:02_related}

\myitparagraph{Automated issue report classification (IRC)} Early IRC studies employed data mining and ML classifiers~\cite{antoniol2008bug, kallis2019ticket, fan2017road}. For example, Antoniol et al.~\cite{antoniol2008bug} leveraged a naive Bayes classifier to distinguish bugs from other issue reports. However, these methods yielded limited performance due to their shallow semantic understanding of issue report text~\cite{heo2025study}. To overcome this limitation, more recent work has adopted BERT-like transformer models for their stronger contextual understanding~\cite{colavito2022issue, bharadwaj2022github, trautsch2022predicting, izadi2022catiss, izadi2022predicting}. These studies typically fine-tune the models on historical issues in a supervised setting. For instance, Izadi et al.~\cite{izadi2022catiss} fine-tuned a RoBERTa~\cite{liu2019roberta} model to classify issues into \textit{Bugs}, \textit{Features}, or \textit{Questions}. Similarly, Trautsch et al.~\cite{trautsch2022predicting} trained a seBERT~\cite{von2022validity} model to classify issues into the same three categories. Although transformer-based approaches achieve strong results, they rely on large amounts of training data, which limits their generalizability to settings where labeled data is scarce~\cite{heo2025study, izadi2022catiss, izadi2022predicting, colavito2022issue}. These challenges have motivated the exploration of LLMs for IRC. Their extensive pretraining allows these models to be applied with little or no task-specific training, reducing the dependence on the large labeled datasets that supervised classifiers require.~\cite{aracena2025applying, aracena2024applyinglargelanguagemodels, colavito2024leveraging, colavito2025benchmarking, colavito2024large}. However, state-of-the-art LLM-based methods do not leverage this property. Instead, they fine-tune the LLM with Low-Rank Adaptation (LoRA~\cite{hu2022lora}) on labeled issue datasets to achieve their best reported performance~\cite{heo2025study, aracena2025applying, aracena2024applyinglargelanguagemodels}. Although fine-tuning is effective, it reintroduces the dependency on labeled training data, remains computationally expensive, and must be repeated whenever the underlying model is swapped or as newly labeled issues accumulate in the project~\cite{fan2017road, kallis2019ticket}.


\myitparagraph{Retrieval-augmented few-shot prompting} Studies show that few-shot prompting can improve LLM performance in specialized tasks due to in-context learning~\cite{assi2026llm, le2023log, ma2023fairnessguidedfewshotpromptinglarge, logan2021cuttingpromptsparameterssimple, brown2020language}. Few-shot prompting has also been applied in IRC. Colavito et al.~\cite{colavito2025benchmarking, colavito2024leveraging} randomly selected one or two labeled examples from their dataset and presented them in a few-shot prompt. Few-shot performance, however, is sensitive to the quality of in-context examples. Research shows that using retrieval techniques to select the most semantically relevant in-context examples can outperform random selection~\cite{liu2022makes, yu2023retrieval}. Since retrieval-augmented few-shot prompting can enhance performance without the overhead of training, it can be a promising alternative to LoRA fine-tuning IRC baseline. To the best of our knowledge, however, this method has not been systematically and empirically evaluated for IRC.
